\documentclass{article}
\usepackage{graphicx} % Required for inserting images
\usepackage{caption} % Required for \captionof
\usepackage{amsmath}
\usepackage{float}
\usepackage{comment}
\usepackage[T1]{fontenc}
\usepackage[utf8]{inputenc}
\usepackage{lmodern}
\usepackage{microtype}
\usepackage{xcolor}
\usepackage{framed}
\renewenvironment{leftbar}{%
  \def\FrameCommand{\textcolor{orange}{\vrule width 3pt} \hspace{10pt}}%
  \MakeFramed {\advance\hsize-\width \FrameRestore}}%
 {\endMakeFramed}

\title{Mechanikai mérés}
\author{Kern Luca, Vincze Csongor}
\date{2026 Április 22}

\begin{document}

\maketitle

\section*{6. Feladat: Határozza meg az eszközben található rugó k rugóállandóját!}

A mágneses csillapítókat eltávolítottuk, és a rugó megnyúlását mértük meg a ráakasztott tömeg függvényében.
Egy súly tömege: $m_{\text{súly}} = 25 \text{g} = 0,025 \text{kg}$

\begin{table}[H]
    \centering
    \begin{tabular}{|c|c|c|}
        \hline
        Tömeg [g] & Megnyúlás [mm] & Rugóállandó [N/m]\\ \hline
          0 & & \\ \hline 
          25 & & \\ \hline
          50 & & \\ \hline 
          75 & & \\ \hline
          100 & & \\ \hline 
          125 & & \\ \hline
        
    \end{tabular}
    \caption{Rugó megnyúlása a ráakasztott tömeg függvényében}
    \label{tab:megnyulas}
\end{table}

Ezek alapján kiszámoltuk a rugóállandókat. $F = -D*x$
A kapott rugóállandók átlaga: $D_{\text{rugó}} = \frac{m}{x} = \frac{0,001 \text{kg}}{0,001 \text{m}} = 1 \frac{\text{N}}{\text{m}}$

\end{document}